\chapter{Chapter 5: Evaluation, Guardrails, and Safety}

\section*{Context}

\textbf{Problem statement:} Build evaluation and guardrail systems that reduce hallucinations, policy violations, and sensitive data leakage while preserving utility and maintaining acceptable latency overhead.

\subsection*{Real-World Examples}

\begin{itemize}
  \item A banking assistant passes offline accuracy tests but fails red-team jailbreak testing in staging; layered guardrails with input/output classifiers catch policy bypass attempts.
  \item A healthcare chatbot overblocks harmless prompts until threshold calibration is tuned with risk-weighted evaluation sets, improving utility by 35\% without increasing policy violations.
\end{itemize}

\bigskip\hrule\bigskip

\section*{Step 1: Understand the Problem and Establish Design Scope}

\subsection*{Simulated Dialogue}

\textbf{Interviewer:} Design the evaluation and safety system for an enterprise LLM platform serving multiple risk domains.\\
\textbf{Candidate:} I need to understand the risk tiers, regulatory requirements, and acceptable overhead before designing guardrail depth.

\textbf{Candidate:} What risk domains does this platform serve?\\
\textbf{Interviewer:} Low-risk (internal Q\&A), medium-risk (customer-facing support), and high-risk (financial advice, healthcare triage).

\textbf{Candidate:} What is the acceptable latency overhead for safety checks?\\
\textbf{Interviewer:} $<$ 100ms for input checks, $<$ 150ms for output checks. Must not inflate p95 beyond 200ms total overhead.

\textbf{Candidate:} What evaluation frequency and coverage is expected?\\
\textbf{Interviewer:} Offline eval before every deployment. Online eval (sampled) continuously. Red-team exercises quarterly.

\subsection*{Functional Requirements}

\begin{itemize}
  \item Classify input risk level (prompt injection, PII, policy violation, off-topic) before model invocation.
  \item Validate output for hallucination, unsupported claims, policy compliance, and PII leakage.
  \item Support configurable guardrail depth per risk tier (light for low-risk, heavy for high-risk).
  \item Provide offline evaluation pipeline for model/prompt/retrieval changes with automated quality gates.
  \item Expose safety telemetry: violation rates, block rates, false-positive rates, guardrail latency.
\end{itemize}

\subsection*{Non-Functional Requirements}

\begin{itemize}
  \item \textbf{Latency:} input guard $<$ 100ms p95, output guard $<$ 150ms p95.
  \item \textbf{Coverage:} 100\% of requests pass through input classifier; 100\% of high-risk outputs pass through output verifier.
  \item \textbf{False-positive rate:} $<$ 3\% block rate on benign requests (measured via held-out calibration set).
  \item \textbf{Auditability:} every block, flag, or modification is logged with trace ID and policy version.
  \item \textbf{Freshness:} guardrail model/policy updates deployable within 1 hour without downtime.
\end{itemize}

\subsection*{Back-of-the-Envelope Estimation}

\subsubsection*{Guardrail Throughput}

\begin{itemize}
  \item Platform QPS: 500 (average), 5,000 (peak)
  \item Input classifier: must handle 5,000 QPS at $<$ 100ms $\rightarrow$ requires $\sim$50ms/request on GPU
  \item At 50ms/request with batching: $\sim$20 requests/batch $\rightarrow$ need capacity for 250 batches/sec at peak
  \item GPU requirement for input classifier (small BERT-class model): 2--4 GPUs
  \item Output verifier (for high-risk routes, $\sim$30\% of traffic): 1,500 QPS at $<$ 150ms $\rightarrow$ 1--2 GPUs
\end{itemize}

\subsubsection*{Evaluation Dataset Sizing}

\begin{itemize}
  \item Offline eval set: 2,000--5,000 labeled examples per domain (question, expected answer, retrieved context, policy labels).
  \item Red-team set: 500--1,000 adversarial prompts per quarterly exercise.
  \item Production sampling: 1--5\% of live traffic for online eval (50--250 requests/minute at peak).
\end{itemize}

\subsubsection*{Cost of Guardrails}

\begin{itemize}
  \item Input classifier: $\sim$\$0.001/request (small model inference cost).
  \item Output verifier (LLM-as-judge using smaller model): $\sim$\$0.003/request.
  \item At 500 QPS average: $\sim$\$172/day for full guardrail coverage.
  \item Trade-off: selectively applying output verifier to high-risk routes only reduces cost by $\sim$70\%.
\end{itemize}

\subsubsection*{Latency Budget Allocation}

\begin{itemize}
  \item Total safety overhead target: $<$ 200ms (input + output combined)
  \item Input classifier: 40--80ms (batched GPU inference)
  \item PII scanner: 10--20ms (regex + NER, parallel with classifier)
  \item Output verifier: 80--120ms (smaller LLM call or classifier)
  \item Policy lookup: 5--10ms (in-memory rule engine)
\end{itemize}

\bigskip\hrule\bigskip

\section*{Step 2: Propose High-Level Design and Get Buy-In}

\subsection*{Core Architecture Components}

\begin{itemize}
  \item \textbf{Input Guard:} prompt injection detector, PII scanner, topic classifier, policy enforcer.
  \item \textbf{Risk Tier Router:} selects guardrail depth based on route/domain risk classification.
  \item \textbf{Output Guard:} hallucination/groundedness checker, PII scanner, policy compliance verifier.
  \item \textbf{Eval Pipeline:} offline benchmark runner, online sampled evaluator, red-team test harness.
  \item \textbf{Policy Engine:} versioned policy rules, threshold configs, and guardrail routing tables.
  \item \textbf{Audit Logger:} immutable log of every guardrail decision with trace ID, policy version, and action.
  \item \textbf{Dashboard/Alerting:} violation rate, block rate, false-positive rate, guardrail latency by risk tier.
\end{itemize}

\subsection*{Pipeline A: Guardrail Pipeline (Request-Time)}

\begin{enumerate}
  \item Request arrives; input guard runs in parallel: prompt injection classifier + PII scanner + topic classifier.
  \item Risk tier router determines guardrail depth: light (log-only), medium (flag + soft block), heavy (hard block + human review).
  \item If input guard triggers hard block: return policy violation response, log event, skip model invocation.
  \item If input passes: proceed to model invocation with risk metadata attached.
  \item Model generates response (streaming or batch).
  \item Output guard runs: groundedness check against retrieved context, PII scan, policy compliance check.
  \item If output guard triggers: rewrite (remove PII), flag (add disclaimer), or block (return safe fallback).
  \item Response delivered to user; guardrail decisions logged to audit trail.
\end{enumerate}

\subsection*{Pipeline B: Evaluation Pipeline (Offline + Online)}

\begin{enumerate}
  \item \textbf{Offline eval (pre-deployment):}
  \begin{itemize}
    \item Run model/prompt/retrieval changes against benchmark eval set (2K--5K examples).
    \item Compute metrics: Recall@K, citation precision, faithfulness, hallucination rate, policy violation rate.
    \item Quality gate: no regression $>$ 2\% on any metric vs baseline.
    \item If gate passes: promote to canary. If fails: block deployment, report regression details.
  \end{itemize}

  \item \textbf{Online eval (continuous):}
  \begin{itemize}
    \item Sample 1--5\% of production traffic.
    \item Run LLM-as-judge on sampled responses for faithfulness and policy compliance.
    \item Track metrics over rolling 24h window with alerting on degradation.
  \end{itemize}

  \item \textbf{Red-team eval (periodic):}
  \begin{itemize}
    \item Run adversarial prompt suite against production system.
    \item Measure jailbreak success rate, policy bypass rate, PII extraction success.
    \item Results feed guardrail model retraining and policy updates.
  \end{itemize}
\end{enumerate}

\textit{[Insert Architecture Diagram: Request flow through input guard (parallel classifiers) $\rightarrow$ risk tier router $\rightarrow$ model invocation $\rightarrow$ output guard (groundedness + PII + policy) $\rightarrow$ response. Eval pipeline: offline benchmarks $\rightarrow$ quality gate $\rightarrow$ canary. Online: sampled traffic $\rightarrow$ LLM-as-judge $\rightarrow$ dashboard. Audit logger connected to all guardrail decisions.]}

\subsection*{Mermaid Architecture Diagram}

\begin{verbatim}
flowchart LR
  request[Request] --> inputGuard[InputGuard]

  inputGuard --> injectionDetector[PromptInjectionDetector]
  inputGuard --> piiScanner1[PIIScanner]
  inputGuard --> topicClassifier[TopicClassifier]

  injectionDetector --> riskRouter[RiskTierRouter]
  piiScanner1 --> riskRouter
  topicClassifier --> riskRouter

  riskRouter -->|"blocked"| blockResponse[PolicyViolationResponse]
  riskRouter -->|"passed"| modelInvoke[ModelInvocation]

  modelInvoke --> outputGuard[OutputGuard]
  outputGuard --> groundCheck[GroundednessChecker]
  outputGuard --> piiScanner2[OutputPIIScanner]
  outputGuard --> policyCheck[PolicyComplianceChecker]

  groundCheck --> responseBuilder[ResponseBuilder]
  piiScanner2 --> responseBuilder
  policyCheck --> responseBuilder

  responseBuilder --> client[Client]

  auditLog[AuditLogger] --> inputGuard
  auditLog --> outputGuard

  evalPipeline[OfflineEvalPipeline] --> qualityGate[QualityGate]
  onlineEval[OnlineSampledEval] --> dashboard[Dashboard]
  redTeam[RedTeamHarness] --> dashboard
\end{verbatim}

\bigskip\hrule\bigskip

\section*{Step 3: Design Deep Dive}

Focus areas for this chapter:

\begin{enumerate}
  \item \textbf{RAG Evaluation Methodology}
  \item \textbf{Guardrail Pipeline Design}
  \item \textbf{Safety Testing and Red-Teaming}
\end{enumerate}

\subsection*{1) RAG Evaluation Methodology}

\subsubsection*{RAG Metric Trade-off Table}

\small
\begin{tabular}{|p{0.13\linewidth}|p{0.17\linewidth}|p{0.17\linewidth}|p{0.17\linewidth}|p{0.20\linewidth}|}
\hline
\textbf{Metric} & \textbf{What It Measures} & \textbf{Strengths} & \textbf{Weaknesses} & \textbf{Use When} \\
\hline
Recall@K & Retrieved relevant docs in top K & Simple, captures retrieval quality & Ignores generation quality & Retrieval-only eval \\
\hline
Citation precision & Generated claims supported by cited source & Directly measures grounding & Requires claim-level annotation & High-stakes RAG \\
\hline
Faithfulness & Answer consistency with retrieved context & End-to-end quality signal & Expensive (LLM-as-judge or human) & Release gating \\
\hline
Unsupported-claim rate & Fraction of claims not grounded in context & Catches hallucination directly & Binary; misses partial support & Safety-critical domains \\
\hline
Answer relevance & Response addresses user query & Captures user satisfaction proxy & Subjective, hard to automate & User-facing quality \\
\hline
\end{tabular}
\normalsize

\textbf{LLM-as-judge evaluation:}

\begin{itemize}
  \item Use a separate (often larger) model to judge faithfulness and relevance.
  \item Anchor with reference answers and rubric to reduce variance.
  \item Calibrate judge agreement against human labels: target Cohen's kappa $>$ 0.7.
  \item Cost: $\sim$\$0.003--0.01 per judged response (using smaller judge model).
  \item Latency: 500ms--2s per judgment $\rightarrow$ not suitable for real-time; use for offline and sampled online eval.
\end{itemize}

\textbf{Evaluation dataset construction:}

\begin{itemize}
  \item 200 queries per domain, manually curated with: query, ideal answer, required citations, policy-relevant edge cases.
  \item Include adversarial queries: ambiguous, multi-hop, unanswerable, policy-boundary.
  \item Version datasets alongside model/prompt versions for reproducibility.
  \item Refresh quarterly to prevent benchmark overfitting.
\end{itemize}

\subsection*{2) Guardrail Pipeline Design}

\subsubsection*{Guardrail Placement Trade-off Table}

\small
\begin{tabular}{|p{0.14\linewidth}|p{0.16\linewidth}|p{0.16\linewidth}|p{0.16\linewidth}|p{0.22\linewidth}|}
\hline
\textbf{Placement} & \textbf{Latency Overhead} & \textbf{Coverage} & \textbf{Complexity} & \textbf{Best For} \\
\hline
Input-only & 40--80ms & Prevents bad inputs & Misses hallucination & Low-risk routes \\
\hline
Output-only & 80--150ms & Catches generation errors & Bad inputs waste compute & Medium-risk routes \\
\hline
Input + output & 120--230ms & Full coverage & Highest overhead & High-risk routes \\
\hline
Async output (log + review) & 0ms (request path) & Full coverage, delayed action & Violations detected after delivery & Batch/async routes \\
\hline
\end{tabular}
\normalsize

\textbf{Input guard components:}

\begin{itemize}
  \item \textbf{Prompt injection detector:} fine-tuned classifier (BERT-class) detecting instruction override attempts. Precision target: $>$ 95\%. Recall target: $>$ 85\%.
  \item \textbf{PII scanner:} regex patterns (SSN, credit card, email) + NER model for names/addresses. Runs in parallel with classifier.
  \item \textbf{Topic classifier:} off-topic detection to prevent misuse (e.g., coding assistant asked for medical advice).
  \item \textbf{Policy enforcer:} rule engine checking input against versioned policy (blocked topics, required disclaimers).
\end{itemize}

\textbf{Output guard components:}

\begin{itemize}
  \item \textbf{Groundedness checker:} compares generated claims against retrieved context. Can be classifier-based (fast, $\sim$50ms) or LLM-based (thorough, $\sim$500ms). Use classifier for real-time, LLM for sampled audit.
  \item \textbf{Output PII scanner:} catches PII leaked from training data or retrieved documents (same pipeline as input but applied to output).
  \item \textbf{Policy compliance checker:} ensures output format, tone, and content adhere to domain policy.
  \item \textbf{Action on violation:} rewrite (strip PII), flag (add disclaimer), soft block (ask user to rephrase), hard block (return safe fallback).
\end{itemize}

\textbf{Risk-tiered guardrail depth:}

\small
\begin{tabular}{|p{0.18\linewidth}|p{0.22\linewidth}|p{0.26\linewidth}|p{0.22\linewidth}|}
\hline
\textbf{Risk Tier} & \textbf{Input Guard} & \textbf{Output Guard} & \textbf{Eval Frequency} \\
\hline
Low (internal Q\&A) & Injection detector only & Log-only (async) & Monthly offline \\
\hline
Medium (customer support) & Full input suite & Classifier-based groundedness & Weekly offline, daily sampled \\
\hline
High (finance, healthcare) & Full input suite + human-in-loop escalation & LLM-based groundedness + PII + policy & Daily offline, continuous sampled \\
\hline
\end{tabular}
\normalsize

\subsection*{3) Safety Testing and Red-Teaming}

\subsubsection*{Red-Team Testing Framework}

\textbf{Attack categories:}

\begin{itemize}
  \item \textbf{Prompt injection:} direct (``ignore previous instructions'') and indirect (injected via retrieved documents).
  \item \textbf{Jailbreak:} role-play, encoding tricks, multi-turn escalation.
  \item \textbf{PII extraction:} probing for training data memorization.
  \item \textbf{Policy bypass:} edge cases at policy boundaries (e.g., borderline medical advice).
  \item \textbf{Denial of service:} extremely long prompts, rapid-fire requests, resource exhaustion.
\end{itemize}

\textbf{Red-team process:}

\begin{enumerate}
  \item Build adversarial prompt suite (500--1,000 prompts per category).
  \item Run against production-equivalent system (same guardrails, same model).
  \item Measure: jailbreak success rate (target: $<$ 2\%), PII extraction rate (target: 0\%), policy bypass rate (target: $<$ 1\%).
  \item Failures feed guardrail model retraining and policy rule updates.
  \item Repeat quarterly; expand suite based on emerging attack patterns.
\end{enumerate}

\subsubsection*{Safety Evaluation Trade-off Table}

\small
\begin{tabular}{|p{0.16\linewidth}|p{0.18\linewidth}|p{0.12\linewidth}|p{0.14\linewidth}|p{0.24\linewidth}|}
\hline
\textbf{Method} & \textbf{Coverage} & \textbf{Cost} & \textbf{Frequency} & \textbf{Best For} \\
\hline
Rule-based checks & Known patterns only & Very low & Every request & PII, format violations \\
\hline
Classifier-based detection & Learned patterns & Low & Every request & Injection, toxicity \\
\hline
LLM-as-judge & Nuanced policy checks & Medium & Sampled (1--5\%) & Faithfulness, policy compliance \\
\hline
Human review & Highest quality & Very high & Escalated cases only & Ambiguous policy decisions \\
\hline
Red-team exercises & Adversarial robustness & High (labor) & Quarterly & Attack surface discovery \\
\hline
\end{tabular}
\normalsize

\subsection*{Production Bottlenecks and Mitigations}

\begin{itemize}
  \item \textbf{Guardrail latency inflating p95}
  \begin{itemize}
    \item Symptom: safety overhead exceeds 200ms budget on interactive routes.
    \item Mitigation: run input classifiers in parallel, use lightweight models, apply heavy checks only on high-risk routes.
  \end{itemize}

  \item \textbf{False-positive overblocking}
  \begin{itemize}
    \item Symptom: $>$ 3\% of benign requests blocked; user complaints rise.
    \item Mitigation: threshold calibration on held-out benign set, per-domain threshold tuning, appeal/override flow.
  \end{itemize}

  \item \textbf{Classifier threshold drift}
  \begin{itemize}
    \item Symptom: block rate changes without policy changes due to input distribution shift.
    \item Mitigation: track threshold stability metrics, re-calibrate monthly, alert on $>$ 1\% block rate deviation.
  \end{itemize}

  \item \textbf{Eval blind spots from synthetic-only datasets}
  \begin{itemize}
    \item Symptom: pass offline eval but fail on real user queries.
    \item Mitigation: supplement synthetic benchmarks with anonymized production query samples.
  \end{itemize}

  \item \textbf{Policy version mismatch during deployment}
  \begin{itemize}
    \item Symptom: new model deployed with old guardrail config; violations spike.
    \item Mitigation: atomic deployment of model + guardrail config as versioned unit.
  \end{itemize}

  \item \textbf{Audit log volume and query performance}
  \begin{itemize}
    \item Symptom: audit queries slow as log volume grows.
    \item Mitigation: structured logging, time-partitioned storage, indexed by trace ID and policy version.
  \end{itemize}
\end{itemize}

\subsection*{Advanced Techniques}

\begin{itemize}
  \item \textbf{LLM-as-judge with calibrated anchors} for nuanced evaluation beyond binary classification.
  \item \textbf{Policy-as-code} guardrails with version-controlled rules in Git, deployed via CI/CD.
  \item \textbf{Risk-tiered abstention} where model returns ``I don't know'' with confidence score rather than generating potentially harmful content.
  \item \textbf{Constitutional AI principles} as alignment guardrails during fine-tuning.
  \item \textbf{Watermarking} for detecting model-generated content in audit scenarios.
\end{itemize}

\bigskip\hrule\bigskip

\section*{Step 4: Wrap Up}

\subsection*{Final Design Summary}

The design makes safety measurable and tunable: layered guardrails with risk-tiered depth, calibrated thresholds tied to false-positive budgets, and evaluation loops that gate every deployment.

\subsection*{Wrap-Around Concerns}

\begin{itemize}
  \item \textbf{Release engineering:} guardrail config deployed atomically with model version; quality gate blocks promotion on any metric regression.
  \item \textbf{Governance:} immutable audit logs with trace IDs, policy version tracking, quarterly red-team reports.
  \item \textbf{Cost controls:} guardrail compute budget per risk tier; selective application of expensive checks to high-risk routes only.
  \item \textbf{Drift controls:} monthly threshold recalibration, weekly eval set refresh, production sampling for distribution shift detection.
\end{itemize}

\bigskip\hrule\bigskip

\section*{Step 5: Interviewer Curveballs}

\subsection*{Interviewer Follow-ups \& Frequently Asked Questions (FAQ)}

\textbf{Q: Your guardrails are blocking 8\% of benign requests. How do you fix this without reducing safety?}

\textbf{The Ideal Answer:} Analyze blocked requests to identify false-positive patterns. Raise classifier threshold on specific categories where precision is low. Add a ``soft block'' tier that flags instead of blocks.

\textbf{Deep Dive \& Trade-offs:}
\begin{itemize}
  \item Threshold tuning is per-category, not global. Injection detector may need different threshold than toxicity classifier.
  \item Soft block (flag + disclaimer) preserves utility while logging potential issues for review.
  \item Hold-out calibration set of 500+ benign examples ensures threshold changes are validated.
\end{itemize}

\textbf{The Interviewer's Rubric:} Checking nuanced threshold management beyond binary block/allow.

\textbf{Q: How do you evaluate faithfulness when you don't have ground-truth answers?}

\textbf{The Ideal Answer:} Use LLM-as-judge with retrieved context as reference. The judge checks whether each generated claim is supported by the provided context, not whether it matches a golden answer.

\textbf{Deep Dive \& Trade-offs:}
\begin{itemize}
  \item Faithfulness evaluation requires only context + response, not golden answers.
  \item Calibrate judge against 200+ human-labeled examples to measure agreement (target kappa $>$ 0.7).
  \item Supplement with unsupported-claim rate as cheaper proxy metric.
\end{itemize}

\textbf{The Interviewer's Rubric:} Checking evaluation methodology sophistication for RAG systems.

\textbf{Q: A red-team exercise reveals a jailbreak that bypasses your input classifier. What do you do?}

\textbf{The Ideal Answer:} Add the attack pattern to training data, retrain classifier, and add a rule-based catch for the specific pattern as immediate mitigation. Then evaluate if output guard would have caught the harmful response.

\textbf{Deep Dive \& Trade-offs:}
\begin{itemize}
  \item Rule-based patch is fastest (deploy in minutes) but brittle against variations.
  \item Classifier retraining is more robust but takes hours/days.
  \item Layered defense: if input guard misses, output guard should still catch harmful content.
  \item Defense in depth is the principle: no single layer should be a single point of failure.
\end{itemize}

\textbf{The Interviewer's Rubric:} Checking incident response depth and defense-in-depth thinking.

\textbf{Q: How do you prevent LLM-as-judge from being gamed or miscalibrated?}

\textbf{The Ideal Answer:} Anchor with reference rubrics and calibration examples. Use multiple judge prompts and aggregate. Track judge-human agreement over time and retune when it drops below 0.7 kappa.

\textbf{Deep Dive \& Trade-offs:}
\begin{itemize}
  \item Single judge prompt is brittle; multi-prompt aggregation reduces variance.
  \item Periodic human audit of judge decisions (100+ samples/month) catches drift.
  \item Different judge model than production model reduces self-evaluation bias.
\end{itemize}

\textbf{The Interviewer's Rubric:} Checking robustness of evaluation infrastructure.

\textbf{Q: What is the latency budget split between input and output guardrails?}

\textbf{The Ideal Answer:} Input guard: 40--80ms (parallel classifiers). Output guard: 80--150ms (groundedness + PII). Total overhead $<$ 200ms. For low-risk routes, skip output guard to stay under 100ms total.

\textbf{Deep Dive \& Trade-offs:}
\begin{itemize}
  \item Running classifiers in parallel (not sequential) is key to meeting latency budget.
  \item Heavyweight LLM-based groundedness check is async for low-risk, synchronous for high-risk.
  \item Pre-compute and cache common policy lookups to avoid per-request DB hits.
\end{itemize}

\textbf{The Interviewer's Rubric:} Checking latency-aware safety system design.

\textbf{Q: How do you handle PII that appears in retrieved documents (not user input)?}

\textbf{The Ideal Answer:} Apply PII scanner to retrieved context before injecting into prompt. Strip or mask PII at the retrieval layer, not just at input/output boundaries.

\textbf{Deep Dive \& Trade-offs:}
\begin{itemize}
  \item PII in retrieved docs can leak into generated responses even if input was clean.
  \item Scanning at retrieval time adds $\sim$10--20ms but prevents downstream leakage.
  \item For high-risk domains: pre-scan and tag documents at indexing time; filter PII-tagged chunks at retrieval.
\end{itemize}

\textbf{The Interviewer's Rubric:} Checking understanding of PII leakage vectors beyond obvious input/output paths.

\textbf{Q: How often should you refresh your evaluation benchmark set?}

\textbf{The Ideal Answer:} Quarterly for core benchmarks. Supplement continuously with anonymized production query samples to catch distribution shift.

\textbf{Deep Dive \& Trade-offs:}
\begin{itemize}
  \item Stale benchmarks lead to Goodhart's law: optimizing for the benchmark rather than real quality.
  \item Production sampling ensures eval reflects actual user behavior.
  \item Version benchmarks alongside model versions for reproducible comparisons.
\end{itemize}

\textbf{The Interviewer's Rubric:} Checking evaluation lifecycle management maturity.

\textbf{Q: What is the minimum safety stack for an MVP launch vs a regulated domain?}

\textbf{The Ideal Answer:} MVP: input injection classifier + output PII scanner + basic topic filter. Regulated: full input/output guard suite + LLM-based groundedness + human-in-loop escalation + audit logging + quarterly red-team.

\textbf{Deep Dive \& Trade-offs:}
\begin{itemize}
  \item MVP stack: $\sim$50ms overhead, $\sim$\$50/day cost. Catches the worst offenders.
  \item Regulated stack: $\sim$200ms overhead, $\sim$\$500/day cost. Defensible under audit.
  \item Key: design the interface to support adding layers without re-architecting.
\end{itemize}

\textbf{The Interviewer's Rubric:} Checking progressive safety maturity planning.

\section*{End of Chapter}

\begin{itemize}
  \item \textbf{Key takeaway:} Safety is not a binary switch but a tunable system. Risk-tiered guardrails, calibrated thresholds, and evaluation loops make safety measurable and improvable.
  \item \textbf{What to practice:} Design a guardrail pipeline with latency budget, specify which checks are synchronous vs async, and explain threshold calibration.
  \item \textbf{Interview drill:} Given a 5\% false-positive block rate on a high-risk route, walk through your diagnosis and calibration workflow.
\end{itemize}
